% GENERATED FILE. Edit canonical lessons, then run npm run generate:books.
% canonical-source-hash: fnv1a64:65c3c6071cb6609c
% canonical-lessons: GU-C129-gali, GU-C129-bag, GU-C129-darvajo, GU-C129-pen, GU-C129-pensil

\chapter{Lane, Park, Gate, Pen, Pencil}
\label{ch:gu-lane-park-gate-pen-pencil}
\hlchaptermodality{129}

\begin{chapteropening}
\textbf{By the end of this chapter:} \emph{I can say a lane, a park, a gate, a pen and a pencil.}

Complete the last lesson of chapter 129: I can say a lane, a park, a gate, a pen and a pencil.
\end{chapteropening}

\section[\texorpdfstring{\gu{ગલી} (galī)}{ગલી (galī)}]{\gu{ગલી} (galī) — a lane}

\label{lesson:GU-C129-gali}



\begin{quote}
Before the new one: say the Gujarati for a library, then the Gujarati for a town square.
\end{quote}

\subsection*{You'll want to know: \gu{ગલી}}
\textbf{\gu{ગલી}} (\emph{galī}) — \textquotedblleft{}a lane\textquotedblright{}.

\begin{grammarlens}[title={What you've built}]
One more word for this chapter.
\end{grammarlens}

\subsection*{Your turn}
\begin{itemize}
\raggedright
  \item \emph{Say it:} \emph{galī}
  \item \emph{Say it:} \emph{galī}, once more
  \item \emph{Say it:} say \emph{galī}
  \item \emph{Recall:} say the Gujarati for a library, then the Gujarati for a town square, then say \emph{galī} again
\end{itemize}

\begin{tcolorbox}[breakable,colback=teal!4,colframe=teal!35!black,title={Before you move on}]
What does \gu{ગલી} mean? (\textquotedblleft{}A lane\textquotedblright{}.) Say it once more.
\end{tcolorbox}
\section[\texorpdfstring{\gu{બાગ} (bāg)}{બાગ (bāg)}]{\gu{બાગ} (bāg) — a park}

\label{lesson:GU-C129-bag}



\begin{quote}
Before the new one: say the Gujarati for a town square, then the Gujarati for a lane.
\end{quote}

\subsection*{You'll want to know: \gu{બાગ}}
\textbf{\gu{બાગ}} (\emph{bāg}) — \textquotedblleft{}a park\textquotedblright{}.

\begin{grammarlens}[title={What you've built}]
One more word for this chapter.
\end{grammarlens}

\subsection*{Your turn}
\begin{itemize}
\raggedright
  \item \emph{Say it:} \emph{bāg}
  \item \emph{Say it:} \emph{bāg}, once more
  \item \emph{Say it:} say \emph{bāg}
  \item \emph{Recall:} say the Gujarati for a town square, then the Gujarati for a lane, then say \emph{bāg} again
\end{itemize}

\begin{tcolorbox}[breakable,colback=teal!4,colframe=teal!35!black,title={Before you move on}]
What does \gu{બાગ} mean? (\textquotedblleft{}A park\textquotedblright{}.) Say it once more.
\end{tcolorbox}
\section[\texorpdfstring{\gu{દરવાજો} (darvājo)}{દરવાજો (darvājo)}]{\gu{દરવાજો} (darvājo) — a gate, a door}

\label{lesson:GU-C129-darvajo}



\begin{quote}
Before the new one: say the Gujarati for a lane, then the Gujarati for a park.

Say \textquotedblleft{}I think that …\textquotedblright{} (\emph{Vichārũ chhũ ke …})
\end{quote}

\subsection*{You'll want to know: \gu{દરવાજો}}
\textbf{\gu{દરવાજો}} (\emph{darvājo}) — \textquotedblleft{}a gate, a door\textquotedblright{}.

\begin{grammarlens}[title={What you've built}]
One more word for this chapter.
\end{grammarlens}

\subsection*{Your turn}
\begin{itemize}
\raggedright
  \item \emph{Say it:} \emph{darvājo}
  \item \emph{Say it:} \emph{darvājo}, once more
  \item \emph{Say it:} say \emph{darvājo}
  \item \emph{Recall:} say the Gujarati for a lane, then the Gujarati for a park, then say \emph{darvājo} again
\end{itemize}

\begin{tcolorbox}[breakable,colback=teal!4,colframe=teal!35!black,title={Before you move on}]
What does \gu{દરવાજો} mean? (\textquotedblleft{}A gate, a door\textquotedblright{}.) Say it once more.
\end{tcolorbox}
\section[\texorpdfstring{\gu{પેન} (pen)}{પેન (pen)}]{\gu{પેન} (pen) — a pen}

\label{lesson:GU-C129-pen}



\begin{quote}
Before the new one: say the Gujarati for a park, then the Gujarati for a gate.
\end{quote}

\subsection*{You'll want to know: \gu{પેન}}
\textbf{\gu{પેન}} (\emph{pen}) — \textquotedblleft{}a pen\textquotedblright{}.

\begin{grammarlens}[title={What you've built}]
One more word for this chapter.
\end{grammarlens}

\subsection*{Your turn}
\begin{itemize}
\raggedright
  \item \emph{Say it:} \emph{pen}
  \item \emph{Say it:} \emph{pen}, once more
  \item \emph{Say it:} say \emph{pen}
  \item \emph{Recall:} say the Gujarati for a park, then the Gujarati for a gate, then say \emph{pen} again
\end{itemize}

\begin{tcolorbox}[breakable,colback=teal!4,colframe=teal!35!black,title={Before you move on}]
What does \gu{પેન} mean? (\textquotedblleft{}A pen\textquotedblright{}.) Say it once more.
\end{tcolorbox}
\section[\texorpdfstring{\gu{પેન્સિલ} (pensil)}{પેન્સિલ (pensil)}]{\gu{પેન્સિલ} (pensil) — a pencil}

\label{lesson:GU-C129-pensil}



\begin{quote}
Before the new one: say the Gujarati for a gate, then the Gujarati for a pen.
\end{quote}

\subsection*{You'll want to know: \gu{પેન્સિલ}}
\textbf{\gu{પેન્સિલ}} (\emph{pensil}) — \textquotedblleft{}a pencil\textquotedblright{}.

\begin{grammarlens}[title={What you've built}]
One more word for this chapter.
\end{grammarlens}

\subsection*{Your turn}
\begin{itemize}
\raggedright
  \item \emph{Say it:} \emph{pensil}
  \item \emph{Say it:} \emph{pensil}, once more
  \item \emph{Say it:} say \emph{pensil}
  \item \emph{Recall:} say the Gujarati for a gate, then the Gujarati for a pen, then say \emph{pensil} again
\end{itemize}

\begin{tcolorbox}[breakable,colback=teal!4,colframe=teal!35!black,title={Before you move on}]
What does \gu{પેન્સિલ} mean? (\textquotedblleft{}A pencil\textquotedblright{}.) Say it once more.
\end{tcolorbox}
